\documentclass[10pt,a4paper]{amsart}
\usepackage[margin=19mm]{geometry}
\usepackage{amsmath,amssymb,mathtools}
\usepackage[scr=rsfs,cal=euler]{mathalfa}
\usepackage{xfrac}
\usepackage[unicode,hidelinks,bookmarks=false]{hyperref}
\newcommand{\ie}{i.e.\ }
\newcommand{\cf}{cf.\ }
\newcommand{\resp}{respectively\ }
\newcommand{\Subsection}{\subsection}
\providecommand{\nobreakdash}{\nobreak\hbox{-}}
\setlength{\emergencystretch}{2em}
\multlinegap0pt
\usepackage[shortlabels]{enumitem}
\usepackage{amssymb}
\usepackage{dsfont}
\usepackage{mathtools}
\usepackage{tikz}
\usepackage{float}
\usepackage{bm}
\usepackage{xcolor}
\newtheorem{definition}{Definition}
\newtheorem{remark}{Remark}
\newtheorem{theorem}{Theorem}
\newtheorem{lemma}{Lemma}
\newtheorem{proposition}{Proposition}
\newcommand{\D}[1]{\mbox{\rm #1}} 
\newcommand{\dd}{\D{d}}
\title{Non-Convex Global Optimization as an Optimal Stabilization Problem: Convergence Rates}
\author{Yuyang Huang, Dante Kalise, Hicham Kouhkouh}
\date{Source: arXiv:2511.11122v2 (CC BY 4.0)}
\begin{document}
\maketitle

\noindent\emph{Source and licence.} Non-Convex Global Optimization as an Optimal Stabilization Problem: Convergence Rates. Version arXiv:2511.11122v2 (CC BY 4.0), distributed by arXiv under the Creative Commons Attribution 4.0 International licence, http://creativecommons.org/licenses/by/4.0/, as recorded in the arXiv licence metadata for this exact version and shown on the abstract page https://arxiv.org/abs/2511.11122v2. Full licence notice: \texttt{licenses/arxiv-2511.11122-CC-BY-4.0.txt}.\par\smallskip

\noindent\textit{Editorial scope.} Counted unit: the article's proposition u (source0.tex lines 1129--1153). The article's complete original proof of it is delivered with it (source0.tex lines 1176--1258, one \texttt{proof} environment of 83 lines, which the article places away from the statement). No mathematics is written, repaired or re-derived: the statement and the proof are the article's own lines, in the article's own notation, and no retained line is substituted or re-flowed.  Retained uncounted as the source-local context the proof needs: lines 308--316; lines 325--330; lines 432--440; lines 729--736; lines 738--753; lines 1036--1046; lines 1048--1053; lines 1068--1075; lines 1086--1093; lines 1100--1106.  Nothing was omitted from any retained span, so the package carries no omission record.  No retained line contains a citation, so the wrapper carries no bibliography and no bibliography entry.  No retained line contains a figure environment, an image-inclusion command or drawing code, so no image is delivered and the package declares no asset dependency.  Editorial formatting only, in addition: an AMS \texttt{amsart} preamble that restates the article's own declarations and macros used by the retained lines, namely the environments \texttt{definition}, \texttt{remark}, \texttt{theorem}, \texttt{lemma}, \texttt{proposition} on separate counters, together with the standard packages those lines need.  The retained environments print in this document as Definition 1, Remark 1, Theorem 1, Lemma 1, Proposition 1 in order of appearance, whereas the article prints the same items with the numbering of its own full document.  The complete native source of the article as distributed in the arXiv e-print for this exact version is bundled unchanged as \texttt{sources/189097/source0.tex}.
\begin{equation}\label{OCP}
\begin{aligned}
     u_{\lambda}(x) & = 
    \inf\limits_{\alpha(\cdot)\in \mathcal{A}} 
    \quad \mathscr{J}(x,\alpha(\cdot))\\
     \text{subject  to  }\; \dot{y}_{x}^{\alpha}(s) & = \alpha(s),\quad y_{x}^{\alpha}(0)=x\in\mathbb{R}^n,\\
     \text{and the controls } &\alpha(\cdot):[0,\infty) \to B_{M} \text{ are measurable},
\end{aligned}
\end{equation}

\begin{enumerate}[label=\textbf{(A\arabic*)}]
\item\label{f: nice} $f : \mathbb{R}^n\to\mathbb{R}$ is  continuous and bounded that is
    \begin{equation*}
    \exists\;\underline{f},\,\overline{f}\; \text{ such that }\; \underline{f} \leq f(x) \leq \overline{f},\quad \forall\;x\in \mathbb{R}^n, \quad \text{where } \; \underline{f}:= \min\limits_{\mathbb{R}^{n}}f.
    \end{equation*}
\end{enumerate}

\begin{definition}\label{def: quasi}
Let $\varepsilon>0$ %be fixed
and $\alpha_{\varepsilon}(\cdot)$ be an admissible control. The corresponding trajectory $y_{x}^{\alpha_{\varepsilon}}(\cdot)$ is called quasi-optimal if 
\begin{equation*}
     \int_{0}^{\infty}\; \left(\frac{1}{2}|\alpha_{\varepsilon}(s)|^{2} + f(y_{x}^{\alpha_{\varepsilon}}(s))\right)\, e^{-\lambda s}\;\text{d}s = \mathscr{J}(x,\alpha_{\varepsilon}(\cdot)) \leq u_{\lambda}(x)+ \varepsilon.
\end{equation*}
When $\varepsilon$ is fixed, then we shall call it an $\varepsilon$-optimal trajectory. In particular, 
optimal trajectories are those corresponding to $\varepsilon=0$. A slightly more general definition will be later introduced in Assumption \ref{quasi opt cost}.
\end{definition}

\begin{enumerate}[label=\textbf{(C)}]
    \item\label{quasi opt cost} There exist $\varepsilon_{\circ}\geq 0$ and $\eta(\cdot)\in L^{\infty}([0,\infty);\mathbb{R}_{+})$ 
    such that
    \begin{equation*}
       \widetilde{\mathscr{J}}(y_{x}^{\alpha_{\varepsilon}}(t),\alpha_{\varepsilon}(t+\,\cdot\,)) \leq \tilde{u}_{\lambda}(y_{x}^{\alpha_{\varepsilon}}(t)) + \varepsilon(t) \quad \quad \text{ for a.a. } t\geq 0,
    \end{equation*} 
    where $\varepsilon(t) = \eta(t)\,\tilde{u}_{\lambda}(y_{x}^{\alpha_{\varepsilon}}(t)) + \varepsilon_{\circ}$.
\end{enumerate}

\begin{remark} \label{rem: (C)}
This more general framework for quasi-optimal trajectories allows for better flexibility in a wide range of applications. The following two observations are worth noting.
\begin{enumerate}
    \item Assumption \ref{quasi opt cost} is simply
    \begin{equation*}
       \widetilde{\mathscr{J}}(y_{x}^{\alpha_{\varepsilon}}(t),\alpha_{\varepsilon}(t+\,\cdot\,)) \leq \big(1+\eta(t)\big)\tilde{u}_{\lambda}(y_{x}^{\alpha_{\varepsilon}}(t)) + \varepsilon_{\circ},
    \end{equation*}
    which coincides with Definition \ref{def: quasi} in the particular case when $t=0$, $\eta(\cdot)\equiv 0$, and then $\varepsilon(\cdot)\equiv \varepsilon_{\circ}$. If moreover $\varepsilon_{\circ}=0$, then this is an optimal trajectory. 
    \item When $t=0$, it also implies 
    \begin{equation*}
       {\mathscr{J}}(x,\alpha_{\varepsilon}(\,\cdot\,)) \leq {u}_{\lambda}(x) + \tilde{\varepsilon} ,
    \end{equation*} 
    where $\tilde{\varepsilon} := \|\eta\|_{\infty}(\overline{f}-\underline{f})/\lambda + \varepsilon_{\circ} \geq \eta(0)\,\tilde{u}_{\lambda}(x) + \varepsilon_{\circ} = \varepsilon(0)$. Hence, if Assumption \ref{quasi opt cost} is satisfied, then the quasi-optimal trajectory is also quasi-optimal in the sense of Definition \ref{def: quasi}, although with a different $\varepsilon$.
\end{enumerate}

\end{remark}

\begin{enumerate}[label=\textbf{(A\arabic*)}]
\setcounter{enumi}{1}
\item\label{f: has min} $f$ attains the minimum, i.e.,  
  \begin{equation*}
    \mathfrak{M}:=\{x\in\mathbb{R}^{n}\,:\, f(x) = \underline{f}:= \min\limits_{z\in\mathbb{R}^{n}}f(z)\} \ne \emptyset .
\end{equation*}
\item\label{f: stab} For all $\delta>0$, there exists  $\gamma=\gamma(\delta)>0$ such that
    \begin{equation*}
         \inf\{f(x) - \underline{f}\,:\,\text{dist}(x,\mathfrak{M})\,> \delta\} \,>\, \gamma(\delta).
    \end{equation*}
\end{enumerate}

\begin{enumerate}[label=\textbf{(D)}]
    \item\label{QG} There exist $\mathbf{r}>0$  and $0<\mathfrak{c}_{1}\leq \mathfrak{c}_{2}<\infty$  such that
    \begin{equation*}
    \frac{\mathfrak{c}_{1}}{2} \, \operatorname{dist}(x,\mathfrak M)^2 \leq \; f(x) - \underline{f} \; \leq  \frac{\mathfrak{c}_{2}}{2} \, \operatorname{dist}(x,\mathfrak M)^2, \quad\forall\,x \in \left\{x\in\mathbb{R}^{n}: \operatorname{dist}(x,\mathfrak{M})\leq \mathbf{r}\right\}.
    \end{equation*}
\end{enumerate}

\begin{theorem}\label{thm: practical}
Let Assumptions \ref{f: nice}, \ref{f: has min}, and \ref{f: stab} be satisfied. Let $x\in\mathbb{R}^{n}$ be given, and $y_{x}^{\alpha_{\varepsilon}}(\cdot)$ be an $\varepsilon$-optimal trajectory for problem \eqref{OCP} with $\varepsilon\geq 0$ small. 
Then for all $\mathbf{r}>0$, there exist $\lambda>0$ small and $\tau = \tau(\mathbf{r},\varepsilon, \lambda)>0$ such that 
\begin{equation*}
    \operatorname{dist}\left(y_{x}^{\alpha_{\varepsilon}}(s), \mathfrak{M}\right) \leq \mathbf{r}, \quad \quad \forall\, s\in [\tau, \infty).
\end{equation*}
In other words, $y_{x}^{\alpha_{\varepsilon}}(s) \in \left\{z\in \mathbb{R}^{n}: \operatorname{dist}(z,\mathfrak{M})\leq \mathbf{r}\right\}$ for all $s\geq \tau$. 
\end{theorem}

\begin{equation}\label{OCP riccati}
\begin{aligned}
    R_{\mathfrak{c}}(x,\mathscr{O}) := \; 
    \inf\limits_{\alpha(\cdot)} \; & \int_{0}^{\infty}\; \left(\frac{1}{2}|\alpha(s)|^{2} + \frac{\mathfrak{c}}{2}\,\operatorname{dist}(y_{x}^{\alpha}(s),\mathscr{O})^{2}\right)\, e^{-\lambda s}\;\dd s,\\
    & \text{subject  to  }\; \dot{y}_{x}^{\alpha}(s) = \alpha(s),\quad y_{x}^{\alpha}(0)=x\in\mathbb{R}^n,\\
    & \text{and the controls } \alpha(\cdot):[0,\infty) \to B_{M} \text{ are measurable}. 
\end{aligned}
\end{equation}

\begin{lemma}\label{lem: ricc}
Let $\mathscr{O}\subset\mathbb{R}^{n}$ be a non-empty and compact set, and let $\mathfrak{c}>0$. Then the value function corresponding to \eqref{OCP riccati} admits the closed form expression
\begin{equation*}
    R_{\mathfrak{c}}(x,\mathscr{O}) = C\,\operatorname{dist}(x,\mathscr{O})^{2} ,\quad \forall\,x\in\mathbb{R}^{n},
\end{equation*}
where $C = (-\lambda + \sqrt{\lambda^{2} + 4\mathfrak{c}}\,)/4 \,>\,0$.
\end{lemma}

\begin{proposition}\label{prop: sandwich u}
For a given $x\in\mathbb{R}^{n}$, let $y_{x}^{\alpha_{\varepsilon}}(\cdot)$ be a quasi-optimal trajectory for problem \eqref{OCP}. % with $\varepsilon\geq 0$ small. 
Suppose the following: 
\begin{enumerate}[label = -]
    \item Assumptions \ref{f: nice}, \ref{f: has min}, and \ref{f: stab} are satisfied, and $\;\mathfrak{M}$ is compact.
    \item Assumption \ref{quasi opt cost} holds for some $\eta(\cdot)\in L^{\infty}([0,\infty);\mathbb{R}_{+})$ and $\varepsilon_{\circ}\geq 0$.
    \item Assumption \ref{QG} holds for some $\mathbf{r}>0$ and $0<\mathfrak{c}_{1}\leq \mathfrak{c}_{2}<\infty$.
\end{enumerate}
Then there exist $\lambda>0$ small and $\tau = \tau(\mathbf{r},\varepsilon_{\circ}, \lambda, \|\eta\|_{\infty})>0$ such that $\;\forall\,t\geq\tau$
\begin{equation*}
    C_{1}\,\operatorname{dist}(y_{x}^{\alpha_{\varepsilon}}(t),\mathfrak{M})^{2} \; 
    \leq \; \tilde{u}_{\lambda}(y_{x}^{\alpha_{\varepsilon}}(t)) + \varepsilon(t) \; 
    \leq \; C_{2}\,\operatorname{dist}(y_{x}^{\alpha_{\varepsilon}}(t),\mathfrak{M})^{2} + \varepsilon(t)+\varepsilon_{\circ}
\end{equation*}
or equivalently, recalling the definition of $\varepsilon(\cdot)$ from Assumption \ref{quasi opt cost},
\begin{equation*}
    C_{1}\,\operatorname{dist}(y_{x}^{\alpha_{\varepsilon}}(t),\mathfrak{M})^{2} \;
    \leq \; (1+\eta(t))\tilde{u}_{\lambda}(y_{x}^{\alpha_{\varepsilon}}(t)) + \varepsilon_{\circ} \; 
    \leq \; (1+\eta(t))\,C_{2}\,\operatorname{dist}(y_{x}^{\alpha_{\varepsilon}}(t),\mathfrak{M})^{2} + (2+ \eta(t))\varepsilon_{\circ},
\end{equation*}
where 
\begin{equation*}
    C_{i} = (-\lambda + \sqrt{\lambda^{2} + 4\mathfrak{c}_{i}}\,)/4 \quad \quad \text{ for } i=1,2.
\end{equation*}
\end{proposition}

\begin{proof}[Proof of Proposition \ref{prop: sandwich u}]
Let $\mathbf{r}>0$ be fixed, and $0<\mathfrak{c}_{1}\leq \mathfrak{c}_{2}<\infty$ be given as in Assumption \ref{QG}. 
The existence of $\lambda>0$ small and $\tau = \tau(\mathbf{r},\varepsilon_{\circ},\lambda, \|\eta\|_{\infty})>0$  such that the quasi-optimal trajectory satisfies for all $s\geq \tau$,  $\;y_{x}^{\alpha_{\varepsilon}}(s) \in \left\{z\in \mathbb{R}^{n}: \operatorname{dist}(z,\mathfrak{M})\leq \mathbf{r}\right\}$, is guaranteed by Theorem \ref{thm: practical} together with Remark \ref{rem: (C)}. Therefore, recalling $\tilde{f}(\cdot)=f-\underline{f}$ and using Assumption \ref{QG}, one obtains
\begin{equation}\label{sandwich f in proof}
    \frac{\mathfrak{c}_{1}}{2} \, \operatorname{dist}(y_{x}^{\alpha_{\varepsilon}}(s),\mathfrak M)^2 \le \; \tilde{f}(y_{x}^{\alpha_{\varepsilon}}(s))  \; \le  \frac{\mathfrak{c}_{2}}{2} \, \operatorname{dist}(y_{x}^{\alpha_{\varepsilon}}(s),\mathfrak M)^2, \quad\forall\,s\geq \tau.
\end{equation}

Let $t\geq \tau$ be arbitrarily fixed, and define $\xi:= y_{x}^{\alpha_{\varepsilon}}(t)$.  The  first inequality in \eqref{sandwich f in proof} yields 
\begin{equation*}
    \int_{t}^{\infty}\left(\frac{1}{2}|\alpha_{\varepsilon}(s)|^{2} + \frac{\mathfrak{c}_{1}}{2} \, \operatorname{dist}(y_{\xi}^{\alpha_{\varepsilon}}(s),\mathfrak M)^2 \right)e^{-\lambda\,(s-t)}\,\dd s
    \leq \; \int_{t}^{\infty}\left(\frac{1}{2}|\alpha_{\varepsilon}(s)|^{2} + \tilde{f}(y_{\xi}^{\alpha_{\varepsilon}}(s))\right)e^{-\lambda\,(s-t)}\,\dd s.
\end{equation*}

The right-hand side is exactly $\widetilde{\mathscr{J}}(\xi,\alpha_{\varepsilon}(t+\,\cdot\,))$. Using the definition of quasi-optimality as in Assumption \ref{quasi opt cost}, the right-hand side term in the above inequality is upper-bounded by $\tilde{u}_{\lambda}(\xi) + \varepsilon(t)$. 

For the left-hand side, it suffices to note that $\alpha_{\varepsilon}(\cdot)$ is admissible (without necessarily being optimal) in the problem \eqref{OCP riccati} where we chose $\mathfrak{c}=\mathfrak{c}_{1}$ and $\mathscr{O}=\mathfrak{M}$. Therefore it is lower bounded by $R_{\mathfrak{c}_{1}}(\xi, \mathfrak{M})$. Ultimately, one gets
\begin{equation}\label{eq: low R}
    R_{\mathfrak{c}_{1}}(\xi, \mathfrak{M}) \leq \tilde{u}_{\lambda}(\xi) + \varepsilon(t).
\end{equation}

On the other hand, let us choose $\beta_{\varepsilon}(\cdot)$ an  $\varepsilon_{\circ}$-optimal control (as in Definition \ref{def: quasi}) in the problem \eqref{OCP riccati} where we chose $\mathfrak{c}=\mathfrak{c}_{2}$ and $\mathscr{O}=\mathfrak{M}$, that is
\begin{equation*}
    \int_{t}^{\infty}\; \left(\frac{1}{2}|\beta_{\varepsilon}(s)|^{2} + \frac{\mathfrak{c}_{2}}{2}\,\operatorname{dist}(y_{\xi}^{\beta_{\varepsilon}}(s),\mathfrak{M})\right)\, e^{-\lambda (s-t)}\;\dd s \leq R_{\mathfrak{c}_{2}}(\xi, \mathfrak{M}) + \varepsilon_{\circ}.
\end{equation*}
Using now the second inequality in \eqref{sandwich f in proof}, one has
\begin{equation*}
    \int_{t}^{\infty}\; \left(\frac{1}{2}|\beta_{\varepsilon}(s)|^{2} + \tilde{f}(y_{\xi}^{\beta_{\varepsilon}}(s))\right)\, e^{-\lambda (s-t)}\;\dd s
    \leq 
    \int_{t}^{\infty}\; \left(\frac{1}{2}|\beta_{\varepsilon}(s)|^{2} + \frac{\mathfrak{c}_{2}}{2}\,\operatorname{dist}(y_{\xi}^{\beta_{\varepsilon}}(s),\mathfrak{M})\right)\, e^{-\lambda (s-t)}\;\dd s.
\end{equation*}
The control $\beta_{\varepsilon}(\cdot)$ being admissible (without necessarily begin optimal) for the problem \eqref{OCP}, one gets
\begin{equation*}
    \tilde{u}_{\lambda}(\xi)
    \leq 
    \int_{t}^{\infty}\; \left(\frac{1}{2}|\beta_{\varepsilon}(s)|^{2} + \tilde{f}(y_{\xi}^{\beta_{\varepsilon}}(s))\right)\, e^{-\lambda (s-t)}\;\dd s
\end{equation*}
and then
\begin{equation}\label{eq: up R}
    \tilde{u}_{\lambda}(\xi)
    \leq 
    R_{\mathfrak{c}_{2}}(\xi, \mathfrak{M}) + \varepsilon_{\circ}.
\end{equation}

Finally, from \eqref{eq: low R} and \eqref{eq: up R} one obtains
\begin{equation*}
\begin{aligned}
    R_{\mathfrak{c}_{1}}(\xi, \mathfrak{M}) 
    & \leq \tilde{u}_{\lambda}(\xi) + \varepsilon(t) 
    = (1+\eta(t))\tilde{u}_{\lambda}(\xi) + \varepsilon_{\circ}\\
    & \leq R_{\mathfrak{c}_{2}}(\xi, \mathfrak{M}) + \varepsilon(t)+\varepsilon_{\circ} \\
    & = R_{\mathfrak{c}_{2}}(\xi, \mathfrak{M})+ \eta(t)\tilde{u}_{\lambda}(\xi) +2\varepsilon_{\circ} \\
    & \leq  (1+\eta(t))\,R_{\mathfrak{c}_{2}}(\xi, \mathfrak{M}) + (2+ \eta(t))\varepsilon_{\circ},
\end{aligned}
\end{equation*}
that is
\begin{equation*}
    R_{\mathfrak{c}_{1}}(\xi, \mathfrak{M}) 
    \leq \tilde{u}_{\lambda}(\xi) + \varepsilon(t) 
    \leq R_{\mathfrak{c}_{2}}(\xi, \mathfrak{M}) + \varepsilon(t)+\varepsilon_{\circ},
\end{equation*}
or, recalling the definition of $\varepsilon(t)$ in Assumption \ref{quasi opt cost}, 
\begin{equation*}
    R_{\mathfrak{c}_{1}}(\xi, \mathfrak{M}) 
    \leq (1+\eta(t))\tilde{u}_{\lambda}(\xi) + \varepsilon_{\circ} 
    \leq (1+\eta(t))\,R_{\mathfrak{c}_{2}}(\xi, \mathfrak{M}) + (2+ \eta(t))\varepsilon_{\circ}.
\end{equation*}
It suffices now to observe, using Lemma \ref{lem: ricc}, that
\begin{equation*}
    R_{\mathfrak{c}_{i}}(\xi,\mathfrak{M}) = C_{i}\,\operatorname{dist}(y_{x}^{\alpha_{\varepsilon}}(t),\mathfrak{M})^{2} \quad \text{ where }  C_{i} = (-\lambda + \sqrt{\lambda^{2} + 4\mathfrak{c}_{i}}\,)/4, \text{ and } i=1,2,
\end{equation*}
which leads to
\begin{equation*}
    C_{1}\,\operatorname{dist}(y_{x}^{\alpha_{\varepsilon}}(t),\mathfrak{M})^{2} 
    \leq \tilde{u}_{\lambda}(\xi) + \varepsilon(t) 
    \leq C_{2}\,\operatorname{dist}(y_{x}^{\alpha_{\varepsilon}}(t),\mathfrak{M})^{2} + \varepsilon(t)+\varepsilon_{\circ},
\end{equation*}
or equivalently
\begin{equation*}
    C_{1}\,\operatorname{dist}(y_{x}^{\alpha_{\varepsilon}}(t),\mathfrak{M})^{2}
    \leq (1+\eta(t))\tilde{u}_{\lambda}(\xi) + \varepsilon_{\circ} 
    \leq (1+\eta(t))\,C_{2}\,\operatorname{dist}(y_{x}^{\alpha_{\varepsilon}}(t),\mathfrak{M})^{2} + (2+ \eta(t))\varepsilon_{\circ}.
\end{equation*}
\end{proof}

\end{document}
